% ============================================================================
% Custom environments preserved from the original AJbook class
% Adapted for use with the ElegantBook template
% ============================================================================

% ============================================================================
% Dynamic watermark (opt-in mechanism)
% Enable from the document preamble with
%   \SeriesWatermarkSetup{<prefix>}{<count>}
% where images <prefix>1.png ... <prefix><count>.png exist (raster images:
% inherently non-selectable in PDF viewers). Placement is deterministic per
% page — count (5..12), golden-angle positions and radii derive from the
% page number, image picked round-robin. Defining \NoWatermark before this
% point disables the mechanism (\SeriesWatermarkSetup becomes a no-op).
% ============================================================================
\RequirePackage{eso-pic}
\RequirePackage{graphicx}
\RequirePackage{tikz}
\RequirePackage{xfp}
\makeatletter

\ifdefined\NoWatermark
  % Watermark disabled.
  \newcommand{\SeriesWatermarkSetup}[2]{%
    \typeout{series-skin: \string\NoWatermark defined; watermark setup ignored}}
\else
\newcommand{\SeriesWatermarkSetup}[2]{%
  \gdef\series@wm@prefix{#1}%
  \gdef\series@wm@n{#2}%
  \AddToShipoutPictureBG{\SeriesDynamicWatermarks}%
}
% Dynamic watermark formula (page-dependent, deterministic + pseudo-random):
%   N(p)   = (p*7+3) mod 8 + 5               — count, 5 to 12
%   θ(p,i) = (p*137 + i*72) mod 360          — angle in degrees (golden angle)
%   r(p,i) = ((p*13 + i*17) mod 7) * 1.5 + 2 — radius 2..11 cm from center
%   pos    = (r*cos θ, r*sin θ) relative to page center
%   image  = <prefix>((p+i) mod count + 1).png
\DeclareRobustCommand{\SeriesDynamicWatermarks}{%
  \begin{tikzpicture}[remember picture, overlay]
    \edef\ajpage{\arabic{page}}%
    \edef\ajnwm{\fpeval{round(\ajpage*7+3-8*floor((\ajpage*7+3)/8)+5)}}%
    \foreach \aji in {1,...,\ajnwm} {%
      \edef\ajangle{\fpeval{\ajpage*137+(\aji-1)*72-360*floor((\ajpage*137+(\aji-1)*72)/360)}}%
      \edef\ajradius{\fpeval{(\ajpage*13+(\aji-1)*17-7*floor((\ajpage*13+(\aji-1)*17)/7))*1.5+2}}%
      \edef\ajxshift{\fpeval{\ajradius*cosd(\ajangle)}}%
      \edef\ajyshift{\fpeval{\ajradius*sind(\ajangle)}}%
      \edef\ajimgidx{\fpeval{round(\ajpage+\aji-1-\series@wm@n*floor((\ajpage+\aji-1)/\series@wm@n)+1)}}%
      \edef\ajwmfile{\series@wm@prefix\ajimgidx.png}%
      \node[anchor=center, opacity=0.10]
        at ([xshift=\ajxshift cm, yshift=\ajyshift cm]current page.center)
        {\includegraphics[width=0.18\paperwidth]{\ajwmfile}};%
    }%
  \end{tikzpicture}%
}
\fi
\makeatother

% ============================================================================
% 选择题选项统一环境
% ============================================================================
% 替代正文里手写的 \vspace{0.35em}\noindent\textbf{A.} 选项块；
% 全书选项的缩进与行距统一由这里控制，正文不再出现逐题手调间距。
\newlist{choices}{enumerate}{1}
\setlist[choices]{
  label = \textbf{\Alph*.},
  labelsep = 0.5em,
  leftmargin = *,
  itemsep = 0.35em,
  topsep = 0.4em,
  parsep = 0pt,
  partopsep = 0pt
}

% ============================================================================
% Redefine example, remark, exercise environments to support list content
% (amsthm-based, compatible with original project usage)
% Task 7: remark/hint sidenote culture; quiet ink chrome; exercise pencil → InkBrown
% ============================================================================
\let\example\relax
\let\endexample\relax
\let\remark\relax
\let\endremark\relax
\let\exercise\relax
\let\endexercise\relax

% Example environment with independent counter (preserving original behavior)
\newtheoremstyle{examplestyle}{3pt}{3pt}{\normalfont}{0pt}{\bfseries\color{main}}{}{0.5em}{\thmname{#1}~\thmnumber{#2}~\thmnote{(#3)}}
\theoremstyle{examplestyle}
\newtheorem{example}{\examplename}[chapter]

% --- Remark: short → margin sidenote; long → quiet main (remark*) ---
% 25mm 边注栏在 footnotesize 下每行约 8 个汉字。
% 超过 64 个 token（约八行）、多段或含展示环境的注释放回正文栏。
% token 计数采用保守上限；remark* 可显式指定正文形式。
\newcounter{remark}[chapter]
\renewcommand{\theremark}{\thechapter.\arabic{remark}}

\ExplSyntaxOn
\bool_new:N \l_series_remark_long_bool
\int_new:N \l_series_remark_nl_int

\str_new:N \l_series_remark_body_str

\cs_new_protected:Npn \series_remark_classify:n #1
  {
    % +b body is tokenized; length ≈ CJK char count via token count.
    % Structure markers via detokenized string (space after csnames).
    \bool_set_false:N \l_series_remark_long_bool
    \int_set:Nn \l_series_remark_nl_int { \tl_count:n {#1} }
    \int_compare:nNnT { \l_series_remark_nl_int } > { 64 }
      { \bool_set_true:N \l_series_remark_long_bool }
    \bool_if:NF \l_series_remark_long_bool
      {
        \str_set:Nx \l_series_remark_body_str { \tl_to_str:n {#1} }
        \exp_args:NnV \regex_match:nnTF
          {
            \\begin\ \{
            ( equation | align | alignat | flalign | gather | multline |
              itemize | enumerate | compactitem | compactenum |
              tikzpicture | multicols | choices | figure | table | tabular | tabularx )
            \*?
            \}
            |
            \\\[
            |
            \\par
          }
          \l_series_remark_body_str
          { \bool_set_true:N \l_series_remark_long_bool }
          { }
      }
  }

\cs_new_protected:Npn \series_remark_head:n #1
  {
    \textbf
      {
        \remarkname~\theremark
        \tl_if_novalue:nF {#1} { \ (#1) }
      }
  }

\cs_new_protected:Npn \series_remark_margin:nn #1 #2
  {
    \refstepcounter{remark}
    \sidenote
      {
        \series_remark_head:n {#1}
        \par
        #2
      }
  }

\cs_new_protected:Npn \series_remark_main:nn #1 #2
  {
    \refstepcounter{remark}
    \begin{tcolorbox}[series-gray-frame]
    \use:c { series@makefloatslocal }
    {
      \footnotesize
      \color{InkMuted}
      \series_remark_head:n {#1}
    }
    \par\nobreak
    {
      \footnotesize
      \color{InkMuted}
      \raggedright
      #2
      \par
    }
    \end{tcolorbox}
  }

\cs_new_protected:Npn \series_remark_auto:nn #1 #2
  {
    \series_remark_classify:n {#2}
    \bool_if:NTF \l_series_remark_long_bool
      { \series_remark_main:nn {#1} {#2} }
      { \series_remark_margin:nn {#1} {#2} }
  }

% Document-level aliases (underscore names only valid under ExplSyntaxOn)
\cs_set_eq:NN \SeriesRemarkAuto \series_remark_auto:nn
\cs_set_eq:NN \SeriesRemarkMain \series_remark_main:nn
\ExplSyntaxOff

\NewDocumentEnvironment{remark}{o +b}{%
  \SeriesRemarkAuto{#1}{#2}%
}{}

% Force main-column quiet form (long remarks / manual override)
\NewDocumentEnvironment{remark*}{o +b}{%
  \SeriesRemarkMain{#1}{#2}%
}{}

% ============================================================================
% Unified exercise counter (shared by exercise, exercises, Exercises)
% ============================================================================
\newcounter{bookexercise}[chapter]
\renewcommand{\thebookexercise}{\thechapter.\arabic{bookexercise}}
\renewcommand{\theHbookexercise}{\thechapter.\arabic{bookexercise}}

% Per-chapter counter for exsol blocks, kept in sync with bookexercise
% (1:1 correspondence between exercises and exsol blocks within each chapter).
\newcounter{exsolcount}[chapter]

% Exercise environment (supporting lists, with pencil icon)
\newtheoremstyle{exercisestyle}{3pt}{3pt}{\normalfont}{0pt}{\bfseries\color{main}}{}{0.5em}{\makebox[-3pt][r]{\scriptsize\color{InkBrown}\ding{46}\quad}\thmname{#1}~\thmnumber{#2}~\thmnote{#3}}
\theoremstyle{exercisestyle}
\newtheorem{exercise}[bookexercise]{\exercisename}

% ============================================================================
% Proof environment with end-of-proof symbol
% Optional note renders in parentheses: \begin{proof}[柯西归纳法] → 证明 (柯西归纳法)
% (elegantbook.cls defines proof without an optional argument, so a [note]
%  leaked into the body as literal square brackets. Consuming it here and
%  wrapping it in parentheses matches the amsthm \thmnote{(#3)} convention.)
% ============================================================================
\let\proof\relax
\let\endproof\relax

\NewDocumentEnvironment{proof}{o}{%
    \par\vspace{0.5em}\noindent
    \textbf{\color{second}\proofname\IfValueT{#1}{ (#1)}\;}%
    \color{black!90}\cfs
}{%
    \hfill$\square$\quad
    \par\vspace{0.5em}%
}

% ============================================================================
% Property environment: accept optional note in parentheses
% (elegantbook.cls defines property without an optional argument; same leak as
%  proof. Override here — loaded after the class, so this wins over both the
%  color-mode and monochrome-mode definitions in the class.)
% Wrapped in the gray solid frame inline (xparse environments cannot be
%  re-wrapped via \let + \RenewDocumentEnvironment — that recurses through
%  xparse's shared internals).
% ============================================================================
\let\property\relax
\let\endproperty\relax
\NewDocumentEnvironment{property}{o}{%
    \csname series@thmcheck\endcsname
    \begin{tcolorbox}[series-gray-frame]
    \csname series@makefloatslocal\endcsname
    \par\noindent\textbf{\color{black}\propertyname\IfValueT{#1}{ (#1)}} \citshape
}{%
    \par
    \end{tcolorbox}%
    \csname series@thmmark\endcsname{property}%
}

% ============================================================================
% Exercise environments (including two-column layouts)
% ============================================================================

% Lower-case exercises environment for per-section exercises.
% Uses the shared bookexercise counter (per-chapter, continues across blocks).
\newenvironment{exercises}{%
    \par\vspace{0.5em}%
}{%
    \par\vspace{0.5em}%
}

% Stack two exercises (single-column). Side-by-side minipages at
% 0.47\textwidth are too narrow under series textwidth (~341pt / ~12cm).
\newcommand{\exerciseentry}[2]{%
    \exercisesingle{#1}%
    \exercisesingle{#2}%
}

% Single-column exercise line (also used for odd leftovers).
\newcommand{\exercisesingle}[1]{%
    \par\vspace{0.8em}\noindent
    \refstepcounter{bookexercise}%
    \textbf{\thebookexercise.}\quad #1
    \par
}

% ============================================================================
% Original capitalized Exercises environment for chapter-end sets
% Uses the shared bookexercise counter via \list + \usecounter.
% ============================================================================
\newenvironment{Exercises}{%
    \markright{习题}
    \addcontentsline{toc}{section}{习题}
    \vspace{1em}\begin{center}
        \Large\heiti 习题
    \end{center}\vspace{0.7em}
    \begin{small}%
    \list{\textbf{\thebookexercise.}}{%
        \usecounter{bookexercise}%
        \renewcommand{\theHbookexercise}{Exercise.\thechapter.\arabic{bookexercise}}%
    }%
}{%
    \endlist\end{small}%
}

% ============================================================================
% Exercise solution collection (exsol): captures body, defers to appendix
% exsol replaces solution for exercise contexts; solution stays for examples.
% ============================================================================
\makeatletter
% Flag: true in noanswer mode (solutions not collected, appendix empty)
\newif\ifseries@noanswer
\AtBeginDocument{%
  \@ifundefined{ELEGANT@result}{}{%
    \ifdefstring{\ELEGANT@result}{noanswer}{%
      \series@noanswertrue
    }{}%
  }%
}

% Capture chapter title for appendix section headers.
% Only in mainmatter — backmatter chapters must not overwrite.
\let\series@sol@orig@makechapterhead\@makechapterhead
\def\@makechapterhead#1{%
  \if@mainmatter
    \expandafter\gdef\csname series@sol@chtitle@\the\c@chapter\endcsname{#1}%
  \fi
  \series@sol@orig@makechapterhead{#1}%
}

% Track which chapters have solutions (comma-separated list, built globally)
\def\series@sol@chlist{}

% Append solution body to per-chapter storage macro.
% #1 = chapter number, #2 = prefix (already expanded by caller via \xdef),
% #3 = solution body (stored unexpanded)
\newtoks\series@sol@existing
\newtoks\series@sol@prefixtoks
\newtoks\series@sol@bodytoks
\newcommand{\series@sol@append}[3]{%
  \@ifundefined{series@sol@ch@#1}{%
    \expandafter\gdef\csname series@sol@ch@#1\endcsname{}%
    % First exsol for this chapter — add to list
    \ifx\series@sol@chlist\empty
      \xdef\series@sol@chlist{#1}%
    \else
      \xdef\series@sol@chlist{\series@sol@chlist,#1}%
    \fi
  }{}%
  \begingroup
    \series@sol@existing\expandafter\expandafter\expandafter{\csname series@sol@ch@#1\endcsname}%
    \series@sol@prefixtoks\expandafter{#2}%
    \series@sol@bodytoks{#3}%
    \expandafter\xdef\csname series@sol@ch@#1\endcsname{%
      \the\series@sol@existing%
      \the\series@sol@prefixtoks%
      \the\series@sol@bodytoks%
    }%
  \endgroup
}

% Emit collected solutions for a chapter.
% #1 = chapter number
\newcommand{\series@sol@emit}[1]{%
  \@ifundefined{series@sol@ch@#1}{}{%
    \csname series@sol@ch@#1\endcsname
  }%
}

% exsol environment: captures body, stores in per-chapter macro, no body output.
% Each stored block is prefixed with the exercise number and solution header (解)
% matching the solution environment's visual style.
\NewDocumentEnvironment{exsol}{+b}{%
  \ifseries@noanswer\else
    \refstepcounter{exsolcount}%
    \xdef\series@sol@prefix{%
      \noexpand\par\noexpand\noindent
      \noexpand\textbf{\noexpand\color{main}\the\c@chapter.\the\c@exsolcount.}%
      \noexpand\quad
      \noexpand\textbf{\noexpand\color{main}\noexpand\solutionname}%
      \noexpand\ \noexpand\citshape
    }%
    \series@sol@append{\thechapter}{\series@sol@prefix}{#1\par}%
  \fi
}{}

% Public interface: emit collected exercise solutions for a chapter in the appendix.
% #1 = chapter number (matching \thechapter at the point where exsol was used)
\newcommand{\printexercisesolutions}[1]{\series@sol@emit{#1}}

% Auto-emit all collected exercise solutions, grouped by chapter.
% Creates \chapter{习题解答} only if any solutions exist.
% Works for both full-book (main.tex) and single-volume builds — the chapter
% list and titles are captured automatically from whatever chapters are included.
\newcommand{\printallexercisesolutions}{%
  \ifx\series@sol@chlist\empty
  \else
    \chapter{习题解答}
    \@for\series@sol@curch:=\series@sol@chlist\do{%
      \expandafter\series@sol@emitone\expandafter{\series@sol@curch}%
    }%
  \fi
}
\newcommand{\series@sol@emitone}[1]{%
  \@ifundefined{series@sol@ch@#1}{}{%
    \section*{第#1章 \csname series@sol@chtitle@#1\endcsname}%
    \addcontentsline{toc}{section}{第#1章 \csname series@sol@chtitle@#1\endcsname}%
    \csname series@sol@ch@#1\endcsname
  }%
}
\makeatother

% ============================================================================
% Summary environment (ink-brown field; signals "wrap-up"; optional title)
% Dark background with light text so the box reads as a section ending block.
% ============================================================================
\NewDocumentEnvironment{summary}{o}{%
    \par\vspace{1em}
    \begin{tcolorbox}[
        colback=InkWash,
        colframe=InkBrown,
        coltitle=PaperIvory,
        coltext=InkJoint,
        title={\heiti\IfValueTF{#1}{#1}{本节小结}},
        fonttitle=\bfseries\sffamily,
        boxrule=0.5pt,
        arc=1pt,
        left=6pt,
        right=6pt,
        top=4pt,
        bottom=4pt,
    ]
}{%
    \end{tcolorbox}
    \par\vspace{1em}
}

% ============================================================================
% Chapter-opening study-tip box (quiet ink; optional title text)
% ============================================================================
\newtcolorbox{wenxintishi}[1][阅读提示]{
    breakable,
    enhanced,
    width = \textwidth,
    colback = PaperIvory,
    colbacktitle = PaperIvory,
    colframe = InkMuted,
    boxrule=0.4pt,
    coltitle = InkBrown,
    coltext = InkBrown,
    fonttitle = \sffamily\bfseries,
    attach boxed title to top left = {yshift=-\tcboxedtitleheight/2,  xshift=\tcboxedtitlewidth/4},
    boxed title style = {boxrule=0pt, colframe=PaperIvory, colback=PaperIvory},
    before skip = 0.5cm,
    top = 3mm,
    bottom = 3mm,
    title={#1}
}

% ============================================================================
% Hint: short → margin; long → quiet main (same threshold as remark)
% ============================================================================
\ExplSyntaxOn
\cs_new_protected:Npn \series_hint_auto:n #1
  {
    \series_remark_classify:n {#1}
    \bool_if:NTF \l_series_remark_long_bool
      {
        \par\vspace{0.3em}
        \noindent
        {\footnotesize\color{InkMuted}\textbf{提示}}
        \par\nobreak
        {\footnotesize\color{InkMuted}\raggedright #1\par}
        \par\vspace{0.2em}
      }
      {
        \sidenote{\textbf{提示}\par #1}
      }
  }
\cs_set_eq:NN \SeriesHintAuto \series_hint_auto:n
\ExplSyntaxOff

\NewDocumentEnvironment{hint}{+b}{%
  \SeriesHintAuto{#1}%
}{}

% ============================================================================
% Biography box helper (quiet ink)
% ============================================================================
\newcommand{\biographyBox}[3]{%
    % #1: Person name
    % #2: Birth and death years
    % #3: Short biography text
    \par\vspace{0.5em}
    \begin{tcolorbox}[
        colback=PaperIvory,
        colframe=InkMuted,
        coltitle=InkBrown,
        coltext=InkBrown,
        title={\textbf{#1} \, (#2)},
        fonttitle=\bfseries,
        boxrule=0.4pt,
        arc=1pt,
        left=5pt,
        right=5pt,
        top=3pt,
        bottom=3pt
    ]
        #3
    \end{tcolorbox}
    \par\vspace{0.5em}
}

% ============================================================================
% Epigraph helper (quiet muted ink)
% ============================================================================
\newcommand{\epigraph}[2]{%
    % #1: Quoted text
    % #2: Source or author
    \par\vspace{1em}
    \begin{flushright}
        \begin{minipage}{0.7\textwidth}
            \color{InkMuted}\rmfamily #1
            \par\vspace{0.3em}
            \raggedleft\color{InkMuted} --- #2
        \end{minipage}
    \end{flushright}
    \par\vspace{1em}
}

% ============================================================================
% Analysis-track block: content that needs limits/derivatives (Vol IV)
% Lightweight (no tcolorbox body) so floats/figures remain legal.
% ============================================================================
\NewDocumentEnvironment{analysispath}{O{延伸讨论}}{%
  \par\vspace{0.7em}%
  \noindent{\sffamily\bfseries\small\color{InkBrown} #1}\par\nobreak
  \vspace{0.35em}%
  \ignorespaces
}{%
  \par\vspace{0.25em}%
  \vspace{0.5em}%
}

% ============================================================================
% Hollow black-white rectangular frame for theorem-like environments
% Wraps existing amsthm environments with tcolorbox; preserves counters/labels.
% Adjacency check: warns when two theorem-like environments are back-to-back.
% ============================================================================
\makeatletter

% --- Shared frame style: black outline, white (hollow) interior, sharp corners ---
% Unbreakable: a theorem-like box must stay on a single page. With breakable
% removed, tcolorbox treats the box as one unit; when it does not fit at the
% bottom of the current page the page breaker moves it whole to the next page
% instead of splitting it mid-statement. (The wrapped environments hold short
% statements, not proofs, so a box taller than a full page is not a concern.)
\tcbset{
  series-thm-frame/.style={
    colframe=black, colback=white, coltext=black,
    boxrule=0.8pt, arc=0pt, outer arc=0pt,
    left=10pt, right=10pt, top=8pt, bottom=8pt,
    before skip=10pt, after skip=10pt,
  }
}

% --- Gray solid frame: black!10 fill, subtle edge, sharp corners, breakable ---
% Used for example, property, exercise, and main-column remark/hint.
% Breakable because examples and long remarks can span page breaks.
\tcbset{
  series-gray-frame/.style={
    breakable,
    colframe=black!20, colback=black!10, coltext=black,
    boxrule=0.4pt, arc=0pt, outer arc=0pt,
    left=10pt, right=10pt, top=8pt, bottom=8pt,
    before skip=10pt, after skip=10pt,
  }
}

% --- Adjacency check ---
\newcommand{\series@lastthm}{}
% Clear the flag whenever a real paragraph of text begins (robust against
% \section etc. because the kernel para/begin hook stacks, unlike \everypar).
\AddToHook{para/begin}[series-thm]{\gdef\series@lastthm{}}
\newcommand{\series@thmcheck}{%
  \ifx\series@lastthm\@empty\else
    \PackageWarning{series-thm}{%
      Two theorem-like environments are adjacent\MessageBreak
      ('\series@lastthm' followed immediately by another).\MessageBreak
      Insert at least one sentence of body text between them.}%
  \fi
}
\newcommand{\series@thmmark}[1]{\gdef\series@lastthm{#1}}

% --- Helper: wrap a numbered amsthm environment with a frame ---
\newcommand{\series@wrapthm}[1]{%
  \csletcs{series@old@#1}{#1}%
  \csletcs{series@old@end#1}{end#1}%
  \RenewDocumentEnvironment{#1}{o}{%
    \series@thmcheck
    \begin{tcolorbox}[series-thm-frame]
      \color{black}%
      \IfNoValueTF{##1}{\csname series@old@#1\endcsname}{\csname series@old@#1\endcsname[##1]}%
  }{%
      \csname series@old@end#1\endcsname
    \end{tcolorbox}%
    \color{black}%
    \series@thmmark{#1}%
  }%
}
% --- Helper: wrap a starred (unnumbered) amsthm environment ---
\newcommand{\series@wrapthmstar}[1]{%
  \csletcs{series@old@#1*}{#1*}%
  \csletcs{series@old@end#1*}{end#1*}%
  \RenewDocumentEnvironment{#1*}{o}{%
    \series@thmcheck
    \begin{tcolorbox}[series-thm-frame]
      \color{black}%
      \IfNoValueTF{##1}{\csname series@old@#1*\endcsname}{\csname series@old@#1*\endcsname[##1]}%
  }{%
      \csname series@old@end#1*\endcsname
    \end{tcolorbox}%
    \color{black}%
    \series@thmmark{#1*}%
  }%
}

% --- Helper: wrap an environment with a gray solid frame (breakable) ---
% Floats are illegal inside tcolorbox ("Not in outer par mode"), so
% figure/table are locally forced to non-floating [H] placement here.
% \sidenote (=\marginpar) also cannot be placed inside a box — a deferred
% marginpar would trigger "Float(s) lost" at the next \clearpage — so it
% is locally rendered as inline muted text instead.
% (float package already loaded in series-skin/packages.tex.)
\newcommand{\series@makefloatslocal}{%
  \let\series@saved@figure\figure
  \let\series@saved@table\table
  \def\figure{\@ifnextchar[{\series@float@opt@fig}{\series@saved@figure[H]}}%
  \def\series@float@opt@fig[##1]{\series@saved@figure[H]}%
  \def\table{\@ifnextchar[{\series@float@opt@tbl}{\series@saved@table[H]}}%
  \def\series@float@opt@tbl[##1]{\series@saved@table[H]}%
  \def\sidenote{\@ifnextchar[{\series@sn@opt}{\series@sn@plain}}%
  \def\series@sn@opt[##1]##2{\series@sn@body{##2}}%
  \def\series@sn@plain##1{\series@sn@body{##1}}%
}
\newcommand{\series@sn@body}[1]{{\footnotesize\color{InkMuted}（#1）}}

\newcommand{\series@wrapthmgray}[1]{%
  \csletcs{series@old@#1}{#1}%
  \csletcs{series@old@end#1}{end#1}%
  \RenewDocumentEnvironment{#1}{o}{%
    \series@thmcheck
    \begin{tcolorbox}[series-gray-frame]
      \color{black}%
      \series@makefloatslocal
      \IfNoValueTF{##1}{\csname series@old@#1\endcsname}{\csname series@old@#1\endcsname[##1]}%
  }{%
      \csname series@old@end#1\endcsname
    \end{tcolorbox}%
    \color{black}%
    \series@thmmark{#1}%
  }%
}

\series@wrapthm{theorem}
\series@wrapthm{definition}
\series@wrapthm{lemma}
\series@wrapthm{corollary}
\series@wrapthm{proposition}
\series@wrapthm{axiom}
\series@wrapthm{postulate}

\series@wrapthmstar{theorem}
\series@wrapthmstar{definition}
\series@wrapthmstar{lemma}
\series@wrapthmstar{corollary}
\series@wrapthmstar{proposition}
\series@wrapthmstar{axiom}
\series@wrapthmstar{postulate}

% Gray solid frame for non-theorem environments
% (property is wrapped inline in its own definition — see above)
\series@wrapthmgray{example}
\series@wrapthmgray{exercise}
\makeatother

% ----------------------------------------------------------------------------
% Reader-fingerprint / anti-piracy measures live in the using project. Load
% them there, after this file, if needed.
% ----------------------------------------------------------------------------
